% !TEX root = ../main.tex

%----------------------------------------------------------------------
\subsection{Metrics and statistical analysis}
\label{sec:evaluation}
%----------------------------------------------------------------------

We report MAE, root-mean-square error (RMSE), signed bias, Spearman rank
correlation, and $R^2$.  Host- and impurity-macro MAE give each chemical group
equal weight, independent of its number of structures.  Cross-validation
predictions are concatenated only when test folds are disjoint; model seeds
are averaged within a fold before concatenation.  Low-energy behavior is
assessed by MAE for $E_{\mathrm f}\leq 0$, MAE within the true lowest-energy
decile of each pooled test regime, and recall of that decile by the predicted
lowest-energy decile.  These strata are evaluation summaries and are not used
for architecture or hyperparameter selection.

For the architecture factorial, each effect is evaluated as the difference
between the average metric at positive and negative levels of the associated
orthogonal contrast within each paired repeat.  Ninety-five-percent intervals
are obtained by nonparametric bootstrap over the five repeat-level contrasts.
For model comparisons on aligned out-of-fold predictions, we bootstrap the
per-sample difference in absolute error using 10,000 paired resamples.  An
interval excluding zero is treated as evidence for a directional difference;
otherwise the result is reported as inconclusive rather than equivalent.

Uncertainty evaluation includes empirical interval coverage and mean width at
50\%, 80\%, 90\%, and 95\% nominal coverage, Gaussian negative log likelihood,
the continuous ranked probability score, Spearman correlation between
uncertainty and absolute error, and risk--coverage curves.  The area under the
risk--coverage curve (AURC) is accompanied by its excess over oracle
ordering~\cite{elyaniv2010}.
All calibration choices use only the dedicated calibration partition.

No new first-principles calculations are used as external validation in this
study.  Previously generated quantum-chemistry files in the project archive
are not matched to the canonical IMP2D targets and therefore are excluded from
the evidence supporting predictive accuracy or materials discovery claims.
